% 来源：OEBPS/Text/01_15.xhtml（OPF spine 顺序；仅第一册）
\cleardoublepage
\BookEntry[约翰·葛瑞本《深奥的简洁》]{约翰·葛瑞本《深奥的简洁》}

\BookCover{images/书单02深奥的简洁.jpg}{[英] 约翰·葛瑞本著| 张宪润译| 湖南科技出版社2008年版}

\begin{BookSummaryBox}
\textbf{以最简单的方式介绍地球与生命复杂原理的一本书}
\par
\end{BookSummaryBox}

\subsubsection*{这本书谈什么？}

尽管现代科技飞速发展，我们的生活还是避免不了突如其来的意外事件：地震、股市震荡，以及常常出现失误的天气预报。然而，生命本身也是非常复杂的现象，生活在千变万化的世界中，我们又该如何看待呢？在《深奥的简洁》中，科普大师葛瑞本解释了蝴蝶效应、盖娅理论，以及各种关于复杂与混沌的概念，并向我们说明：复杂，其实很简单，即使是看来完全随机的行为深处，也只是遵从简单的因果规律。
\par

从地球起源，到混沌与复杂如何与宇宙各方面紧密地交织在一起，并以相同的方式掌控了生命与银河的演化，以及地球外是否还有生命存在，这本书对这些问题都一一作了解答。
\par

\subsubsection*{芒格为什么谈起这本书？}

在芒格的即席谈话中，有许多幽默的言辞。关于《深奥的简洁》这本书，他说道：
\par

\begin{quotation}
想要理解万物是非常困难的，但如果你无法理解它，你可以把问题交给一位比你聪明的朋友呀。
\par
\end{quotation}

\subsubsection*{《深奥的简洁》作者约翰·葛瑞本如是说}

“我们在周遭世界——甚至在生命世界中——所见到的复杂行为只是‘从深奥的简洁中所浮现出的复杂表象’。”
\par

“想象有一个智慧，它在任何时间都知道所有控制大自然的力量，同时也知道每一项事物的运动状态。假设这智慧体可以将所有资料加以分析，它将能把宇宙中大大小小物体运动的状态用一个公式描述。对它而言，没有不确定的东西；它的眼中可以清楚地看到未来与过去。”
\par
